\section{Giải thích là gì?}
\label{sec:ch01-giai-thich-la-gi}

Trong cuốn sách này, một giải thích là thông tin được trình bày để con người
hiểu một khía cạnh của hành vi mô hình. Định nghĩa này cố ý chưa hứa rằng thông
tin ấy mô tả đúng cơ chế bên trong. Hứa như vậy quá sớm sẽ trộn lẫn hai việc:
trình bày cho người đọc và mô tả điều mô hình đã làm.

\begin{definition}[\tn{Khả năng diễn giải}{interpretability}]
\label{def:ch01-kha-nang-dien-giai}
Khả năng diễn giải là năng lực trình bày hành vi của một hệ thống học máy bằng
những thuật ngữ con người có thể hiểu và dùng cho một mục đích đã nêu.
\end{definition}

Định nghĩa~\ref{def:ch01-kha-nang-dien-giai} có hai phần đáng giữ lại. Thứ
nhất, người nhận giải thích là một phần của bài toán. Bản đồ nhiệt có thể hữu
ích cho chuyên gia ảnh y khoa nhưng vô dụng trong một quy trình khi người vận
hành phải chọn giữa các phương án hành động. Thứ hai, tính hữu ích không đồng
nghĩa với tính đúng về cơ chế.

Một hướng tiếp cận đặt \tn{minh bạch}{transparency} vào chính mô hình. Với
một mô hình tuyến tính nhỏ, ta có thể lần theo biến đầu vào, tham số và phép
tính dẫn đến một dự đoán. Hướng khác tạo \tn{giải thích hậu nghiệm}{post-hoc
explanation} sau khi mô hình đã được huấn luyện: một danh sách đặc trưng, một
bản đồ nhiệt, hoặc một mô hình thay thế cục bộ. Lipton đặt hai hướng này cạnh
nhau vì chúng trả lời những câu hỏi khác nhau~\autocite{p13mythos}.

Không nên biến sự phân biệt này thành một thứ bậc đơn giản. Mô hình minh bạch
cũng có thể quá lớn để người dùng kiểm tra toàn bộ, còn một giải thích hậu
nghiệm có thể giúp phát hiện một lỗi cụ thể. Điểm cần giữ là nguồn gốc của
lời giải thích. Khi nó chỉ xấp xỉ hành vi của mô hình, ta cần một cách kiểm tra
riêng cho quan hệ giữa xấp xỉ và cơ chế.
